% kap07.tex — Virknings-vinkel-variable
% Hamilton-Jacobi-kompendiet

\chapter{Virknings-vinkel-variable}
\label{kap:virkning-vinkel}

\section{Motivation: periodiske systemer uden fuld løsning}

I kapitel~\ref{kap:hj} løste vi den harmoniske oscillator ved at
integrere den tidsuafhængige Hamilton-Jacobi-ligning helt igennem for
at finde $W(x,E)$, og derefter differentiere med hensyn til $E$ for at
finde bevægelsen. Denne fulde løsning gav os $x(t)$ eksplicit — men den
kræver, at vi rent faktisk kan udføre integralet og derefter
\emph{invertere} resultatet for $x(t)$, hvad der langt fra altid er
muligt i lukket form.

For \emph{periodiske} systemer (og alle vores eksempler til videre er
periodiske eller kvasi-periodiske bevægelser i én frihedsgrad) findes
der en mellemvej, som giver os de fysisk vigtigste størrelser —
navnlig svingningsfrekvenserne — uden at kræve den fulde inversion.
Dette er \emph{virknings-vinkel-variablene}\index{virknings-vinkel-variable}. De er ikke kun et regne\-
teknisk trick: de er den formalisme, hvor den ældre kvanteteoris
kvantiseringsregler (kapitel~\ref{kap:adiabatisk}) tager deres mest
naturlige form.

\section{Klassifikation af periodisk bevægelse i én frihedsgrad}

Betragt et system med én frihedsgrad og tidsuafhængig $H(q,p)$. Banen i
faserummet $(q,p)$ — faseportrættet, jf.\ figur~\ref{fig:faseportraet}
— falder i to kvalitativt forskellige klasser:

\begin{itemize}
  \item \textbf{Libration}\index{libration}: banen er en lukket kurve i faserummet (som
        ellipserne for den harmoniske oscillator). Både $q$ og $p$ er
        periodiske funktioner af tiden med samme periode. Dette
        optræder typisk for bunden bevægelse i en potentialbrønd.
  \item \textbf{Rotation}\index{rotation (periodisk bevægelse)}: $p$ er en periodisk funktion af $q$, men
        $q$ selv vokser ubegrænset (som en vinkel, der løber
        rundt gentagne gange) — det klassiske eksempel er et
        pendul, der svinger hele vejen rundt i stedet for at
        oscillere, eller blot en fri partikel på en cirkel.
\end{itemize}

I begge tilfælde er $p(q)$ (banen i faserummet, ved fast $E$) en
\emph{periodisk} funktion af $q$, med en veldefineret periode i
$q$-retningen (for libration: hele svingningen; for rotation: én
omgang). Det er denne fælles egenskab, virknings-vinkel-formalismen
udnytter.

\section{Definition af virkningsvariablen}

For et system i én frihedsgrad definerer vi \emph{virkningsvariablen}\index{virkningsvariabel}
\begin{equation}
  \boxed{J \equiv \oint p\,dq,}
  \label{eq:def-J}
\end{equation}
hvor integralet tages over \'en fuld periode af bevægelsen: for
libration, hele vejen rundt om den lukkede faserumskurve; for
rotation, over ét fuldt gennemløb af $q$ (fx $0$ til $2\pi$ for en
vinkelvariabel). Da $p=\partial W/\partial q$ langs banen (fra
Hamilton-Jacobi-teorien i kapitel~\ref{kap:hj}), er $J$ funktionen af
energien $E$ alene:
\begin{equation}
  J(E) = \oint \pdv{W}{q}\,dq.
  \label{eq:J-fra-W}
\end{equation}
Geometrisk er $J/(2\pi)$ (eller, i visse konventioner, $J$ selv) netop
\emph{arealet}, som den lukkede faserumskurve indeslutter (for
librationstilfældet) — se figur~\ref{fig:virkningsareal}. Dette er en
direkte generalisering af det areal, vi allerede så for den harmoniske
oscillators ellipser i figur~\ref{fig:faseportraet}.

\begin{figure}[htbp]
  \centering
  \begin{tikzpicture}[scale=1.0]
    \draw[-{Latex}] (-2.8,0) -- (2.8,0) node[right] {$q$};
    \draw[-{Latex}] (0,-2.0) -- (0,2.0) node[above] {$p$};
    \draw[thick,blue!60!black,fill=blue!12] (0,0) ellipse (2.2 and 1.5);
    \node[blue!60!black] at (0,0) {$J=\oint p\,dq$};
    \draw[-{Latex},thick] (2.2*0.71,1.5*0.71) -- (2.2*0.71-0.35,1.5*0.71+0.5);
    \node at (2.1,1.55) {bane};
  \end{tikzpicture}
  \caption{Virkningsvariablen $J=\oint p\,dq$ er (op til en faktor
  $2\pi$) arealet, som banen indeslutter i faserummet — her for
  libration (en lukket bane, som for den harmoniske oscillator).}
  \label{fig:virkningsareal}
\end{figure}

\section{Den konjugerede vinkelvariabel}

Da $J=J(E)$ (fra \eqref{eq:J-fra-W}) er en funktion af $E$ alene, kan
vi — i hvert fald i princippet — invertere den til $E=E(J)$, og
udtrykke Hamiltons karakteristiske funktion $W$ som funktion af $q$ og
$J$ i stedet for $q$ og $E$: $W=W(q,J)$. Vi udfører nu en kanonisk
transformation af type $F_2$ (kapitel~\ref{kap:kanoniske-transf}), med
$W(q,J)$ selv som genererende funktion, og med $J$ i rollen som den
nye impuls:
\begin{equation}
  p = \pdv{W}{q}, \qquad w \equiv \pdv{W}{J},
  \label{eq:def-w}
\end{equation}
hvor vi kalder den nye, konjugerede koordinat $w$ — \emph{vinkel\-
variablen}\index{vinkelvariabel}. Da $H=E=E(J)$ i disse variable ikke afhænger af $w$, giver
Hamiltons ligninger \eqref{eq:hamilton-nye-var} med $K=E(J)$:
\begin{equation}
  \dot J = -\pdv{E}{w} = 0, \qquad \dot w = \pdv{E}{J} \equiv \nu(J),
  \label{eq:hamilton-vinkel}
\end{equation}
hvor $\nu(J)\equiv dE/dJ$ er konstant i tiden (fordi $J$ selv er
konstant). Den anden ligning integreres trivielt:
\begin{equation}
  w(t) = \nu t + w_0.
  \label{eq:w-lineaer}
\end{equation}
Vinklen $w$ vokser altså \emph{lineært} med tiden — al den
tilsyneladende komplicerede periodiske bevægelse i $q$ er blevet
transformeret om til en variabel, der simpelthen løber rundt med
konstant hastighed $\nu$.

\subsection{$\nu$ er svingningsfrekvensen}

Vi viser nu, at $\nu$ netop er systemets \emph{svingningsfrekvens}\index{svingningsfrekvens} — hvilket
retfærdiggør navnet ``vinkelvariabel'' og gør hele konstruktionen
praktisk anvendelig. Betragt ændringen i $w$, når $q$ gennemløber netop
\'en fuld periode (defineret som i diskussionen omkring
\eqref{eq:def-J}):
\begin{equation}
  \Delta w \equiv \oint \pdv{w}{q}\,dq = \oint \pdv{}{q}\left(\pdv{W}{J}\right)dq
  = \dv{}{J}\oint \pdv{W}{q}\,dq = \dv{J}{J}\oint p\,dq\Big/ dJ,
\end{equation}
mere ligefremt: vi ombytter $\partial/\partial J$ og $\oint(\cdots)dq$
(tilladt, da integrationskonturen — én fuld periode i $q$ — ikke selv
afhænger af $J$), og bruger \eqref{eq:J-fra-W}:
\begin{equation}
  \Delta w = \dv{}{J}\oint\pdv{W}{q}\,dq = \dv{J}{J} = 1.
  \label{eq:delta-w-1}
\end{equation}
Vinkelvariablen $w$ vokser altså med præcis $1$ for hver fulde periode
af den oprindelige bevægelse i $q$. Sammenholdt med
\eqref{eq:w-lineaer}, hvor $w$ vokser lineært med hastighed $\nu$,
betyder dette, at perioden $\tau$ af den oprindelige bevægelse opfylder
$\nu\tau=1$, altså
\begin{equation}
  \boxed{\nu = \dv{E}{J} = \frac{1}{\tau}.}
  \label{eq:nu-er-frekvens}
\end{equation}
Dette er den centrale pointe i hele kapitlet: \emph{svingnings\-
frekvensen kan beregnes som $dE/dJ$, uden at vi nogensinde behøver
finde bevægelsen $q(t)$ eksplicit}. Vi skal blot kende $J$ som funktion
af $E$ — et almindeligt (ikke partielt) integral, jf.\
\eqref{eq:def-J} — og differentiere.

\section{Eksempel: den harmoniske oscillator, igen}

\begin{eksempel}
Vi genbesøger den harmoniske oscillator en tredje gang (jf.\
kapitel~\ref{kap:hamilton} og \ref{kap:hj}), nu via virknings-vinkel-
metoden. Faserumsbanen ved energi $E$ er ellipsen
$p^2/(2mE)+x^2/(2E/k)=1$ (fra kapitel~\ref{kap:hj}'s resultat
$x=a\sin(\omega(t+\beta))$, $p=ma\omega\cos(\omega(t+\beta))$, med
$a=\sqrt{2E/k}$). Virkningsvariablen \eqref{eq:def-J} er per
definition arealet af denne ellipse, ganget med $2\pi$ (se
figur~\ref{fig:virkningsareal}): en ellipse med halvakser
$a=\sqrt{2E/k}$ (i $x$-retningen) og $b=\sqrt{2mE}$ (i $p$-retningen)
har areal $\pi a b$, så
\begin{equation}
  J = \oint p\,dx = \pi ab = \pi\sqrt{\frac{2E}{k}}\sqrt{2mE} = 2\pi E\sqrt{\frac{m}{k}} = \frac{2\pi E}{\omega},
\end{equation}
hvor vi brugte $\omega=\sqrt{k/m}$. Løst for $E$:
\begin{equation}
  E(J) = \frac{\omega J}{2\pi} = \omega\left(\frac{J}{2\pi}\right).
\end{equation}
Frekvensen, ifølge \eqref{eq:nu-er-frekvens}, er
\begin{equation}
  \nu = \dv{E}{J} = \frac{\omega}{2\pi},
\end{equation}
netop svingningsfrekvensen $\nu=\omega/(2\pi)$ (i enheder af
svingninger pr.\ tidsenhed; $\omega$ selv er vinkelfrekvensen i
radianer pr.\ tidsenhed) — udledt uden på noget tidspunkt at have
integreret $dW/dx$ eksplicit, i modsætning til den fulde løsning i
kapitel~\ref{kap:hj}. Bemærk også formen $E=\nu J$ (med $\nu=\omega/2\pi$):
dette er den klassiske sammenhæng, der — som vi antyder i næste kapitel
— bliver til Plancks kvantiseringsbetingelse $E=nh\nu$, når $J$
begrænses til heltalsmultipla af Plancks konstant $h$.
\end{eksempel}

\section{Flere frihedsgrader: separable systemer}

Konstruktionen generaliseres til $n$ frihedsgrader, når
Hamilton-Jacobi-ligningen er \emph{fuldt separabel}: dvs.\ når man kan
finde en komplet integral af formen
\begin{equation}
  W(q_1,\dots,q_n,\alpha) = \sum_{i=1}^n W_i(q_i,\alpha),
\end{equation}
hvor hver $W_i$ kun afhænger af sin egen koordinat $q_i$ (og af de
$n$ konstanter $\alpha$ som parametre). Dette er tilfældet for alle
systemer, vi har mødt i kompendiet (fx central-kraft-problemet i
planpolar-koordinater, hvor $r$- og $\phi$-delen separerer, jf.\
opgaverne i kapitel~\ref{kap:lagrange} og \ref{kap:hamilton}).

For hvert sådant separabelt ``led'' defineres en egen virkningsvariabel
\begin{equation}
  J_i \equiv \oint p_i\,dq_i,
\end{equation}
hvor $p_i=\partial W_i/\partial q_i$, og integralet tages over den
periode, der er naturlig for netop koordinaten $q_i$ (som kan være
libration eller rotation, forskelligt for hver $i$). Hver $J_i$ har sin
egen konjugerede vinkelvariabel $w_i$, som vokser lineært med sin egen
frekvens $\nu_i=\partial E/\partial J_i$ — nu en \emph{partiel}
afledt, da $E=E(J_1,\dots,J_n)$ generelt afhænger af alle
virkningsvariablene. Systemets fulde bevægelse er dermed opløst i $n$
uafhængige, lineært voksende vinkler, hver med sin egen frekvens — en
\emph{multiperiodisk}\index{multiperiodisk bevægelse} bevægelse.

\begin{bemaerkning}
Dette er den klassiske forløber for begrebet ``normalmoder'' og for
kvantetalene $(n_1,\dots,n_n)$, der mærker de enkelte frihedsgrader i
et kvantiseret system (fx radial- og vinkelkvantetal for
hydrogenatomet). Vi forfølger denne parallel systematisk i næste
kapitel.
\end{bemaerkning}

\subsection{Det geometriske billede: bevægelsestorusen}

Det er værd at gøre sig det geometriske billede klart, fordi det er
selve grundstenen, som al senere teori om integrable og
ikke-integrable systemer bygger videre på. For \'en frihedsgrad
foregår bevægelsen på en lukket kurve i faserummet (figur~
\ref{fig:virkningsareal}) — topologisk en cirkel, parametriseret af
vinklen $w$. For $n$ frihedsgrader, med $n$ uafhængige
virkningsvariable $J_1,\dots,J_n$ alle konstante, foregår bevægelsen
tilsvarende på en $n$-dimensional \emph{torus}\index{bevægelsestorus} i faserummet: hver
vinkel $w_i\in[0,1)$ (eller $[0,2\pi)$, afhængig af konvention) er en
uafhængig cirkulær koordinat på torusen, og systemets tilstand
$(w_1,\dots,w_n)$ løber rundt på den med konstante hastigheder
$\nu_1,\dots,\nu_n$ — se figur~\ref{fig:torus} for $n=2$.

\begin{figure}[htbp]
  \centering
  \begin{tikzpicture}[scale=1.0]
    % Torus (2D projektion): ydre og indre kontur
    \draw[thick,blue!55!black,fill=blue!6] (0,0) ellipse (3.2 and 1.7);
    \draw[thick,blue!55!black,fill=white] (0,0.15) ellipse (1.5 and 0.75);
    % Antydning af "rørets" runding med et par ekstra tynde ellipser
    \draw[blue!35!black,thin] (0,-0.35) ellipse (2.4 and 1.25);
    \draw[blue!35!black,thin] (0,0.6) ellipse (2.75 and 1.45);
    % w_1: "lang vej rundt" (den store cirkel)
    \draw[-{Latex},orange!70!black,thick] (2.35,0.75) arc[start angle=25,end angle=155,x radius=2.6,y radius=1.4];
    \node[orange!70!black] at (0,2.15) {$w_1$ (lang vej rundt om torusen)};
    % w_2: "kort vej rundt" (om røret) - antydet ved en lille lukket kurve
    \draw[-{Latex},green!45!black,thick] (2.2,-0.15) to[out=100,in=-10] (1.85,0.55) to[out=190,in=80] (1.5,-0.15) to[out=-80,in=170] (1.85,-0.75) to[out=10,in=-100] (2.2,-0.15);
    \node[green!45!black] at (2.85,-1.1) {$w_2$ (kort vej, om ``røret'')};
    \node[blue!55!black] at (0,-2.15) {$n=2$: bevægelsen ligger på en 2-torus};
  \end{tikzpicture}
  \caption{For $n$ separable frihedsgrader ligger bevægelsen på en
  $n$-dimensional torus i faserummet (her $n=2$), med hver
  vinkelvariabel $w_i$ som en uafhængig cirkulær koordinat. Er
  frekvenserne $\nu_1,\dots,\nu_n$ \emph{inkommensurable} (deres
  forhold er irrationale), fylder banen hele torusen tæt op
  (\emph{kvasiperiodisk} bevægelse); er de \emph{kommensurable}
  (rationalt forhold), lukker banen sig selv efter endeligt mange
  omløb (\emph{resonans}).}
  \label{fig:torus}
\end{figure}

Sondringen mellem \emph{kommensurable} og \emph{inkommensurable}
frekvenser, nævnt i figurteksten, er ikke en teknisk fodnote — det er
selve nålepunktet, hvor den pæne, fuldstændigt løselige verden, dette
kapitel beskriver, begynder at gå i stykker: perturberes et
integrabelt system svagt, er det netop de \emph{resonante} tori
(kommensurable frekvensforhold) der er skrøbelige og typisk opløses i
kaotisk bevægelse, mens de tilstrækkeligt \emph{ikke-resonante} tori
overlever (om end let deformerede) — indholdet af
Kolmogorov-Arnold-Moser-sætningen (KAM-sætningen). Dette er
selve udgangspunktet for et separat kompendium om kaos og
ikke-lineær dynamik.

\section{Eksempel: Kepler-problemet — flere frihedsgrader og degeneration}
\label{sec:kepler-action-angle}

\begin{eksempel}
Vi regner nu et fuldstændigt eksempel med $n=2$ frihedsgrader
igennem: en partikel med masse $m$, der bevæger sig i et
Kepler-/Coulomb-potentiale $V(r)=-k/r$ ($k>0$) i planpolar-koordinater
$(r,\phi)$. Hamilton-funktionen er
\begin{equation}
  H = \frac{p_r^2}{2m} + \frac{p_\phi^2}{2mr^2} - \frac{k}{r}.
\end{equation}
$\phi$ er cyklisk (jf.\ kapitel~\ref{kap:lagrange}), så $p_\phi=\ell$
(banemomentet) er en bevægelseskonstant, og
Hamilton-Jacobi-ligningen separerer fuldstændigt: $W=W_r(r)+\ell\phi$.
Vi har dermed to virkningsvariable, én for hver frihedsgrad.

\textbf{Den azimutale virkningsvariabel} er triviel, da $p_\phi=\ell$
er konstant, og $\phi$ er en \emph{rotationskoordinat} (den løber
rundt, $0$ til $2\pi$, for hvert omløb):
\begin{equation}
  J_\phi = \oint p_\phi\,d\phi = 2\pi\ell.
  \label{eq:kepler-Jphi}
\end{equation}

\textbf{Den radiale virkningsvariabel} kræver mere arbejde. For en
bunden bane ($E<0$) svinger $r$ periodisk (\emph{libration}) mellem et
minimum $r_{\min}$ og maksimum $r_{\max}$ (perihel og aphel), hvor
$p_r=\sqrt{2m\bigl(E+k/r\bigr)-\ell^2/r^2}=0$. Virkningsvariablen er
\begin{equation}
  J_r = \oint p_r\,dr = 2\int_{r_{\min}}^{r_{\max}}
    \sqrt{2mE + \frac{2mk}{r} - \frac{\ell^2}{r^2}}\;dr,
  \label{eq:kepler-Jr-integral}
\end{equation}
(faktor $2$, fordi $\oint$ dækker både ud- og indadgående bevægelse).
Dette er et standardintegral (løses fx ved konturintegration i det
komplekse $r$-plan, eller ved substitutionen $r=1/u$ efterfulgt af
kvadratkomplettering — se opgave nedenfor), med resultatet
\begin{equation}
  J_r = -2\pi\ell + 2\pi k\sqrt{\frac{m}{-2E}}.
  \label{eq:kepler-Jr}
\end{equation}
Læg mærke til, at $-2\pi\ell=-J_\phi$ indgår, jf.\
\eqref{eq:kepler-Jphi} — det er ingen tilfældighed: banemomentets
bidrag til den \emph{radiale} bevægelses ``omkreds'' i faserummet
skyldes centrifugalleddet $\ell^2/r^2$ i $p_r$, og forsvinder netop
med $J_\phi$'s eget bidrag, når vi lægger de to sammen.

\textbf{Degeneration: energien afhænger kun af summen $J_r+J_\phi$.}
Læg \eqref{eq:kepler-Jphi} og \eqref{eq:kepler-Jr} sammen:
\begin{equation}
  J \equiv J_r + J_\phi = 2\pi k\sqrt{\frac{m}{-2E}}.
\end{equation}
Løs for $E$:
\begin{equation}
  \boxed{E = -\frac{2\pi^2mk^2}{J^2} = -\frac{2\pi^2mk^2}{(J_r+J_\phi)^2}.}
  \label{eq:kepler-E-af-J}
\end{equation}
Dette er et bemærkelsesværdigt resultat: $E$ afhænger \emph{ikke} af
$J_r$ og $J_\phi$ hver for sig, kun af deres \emph{sum}. Systemet er
\emph{degenereret}\index{degeneration (virknings-vinkel-variable)}: de to frekvenser
\begin{equation}
  \nu_r = \pdv{E}{J_r} = \pdv{E}{J}, \qquad
  \nu_\phi = \pdv{E}{J_\phi} = \pdv{E}{J}
\end{equation}
er \emph{identiske}, $\nu_r=\nu_\phi$ — den radiale og den azimutale
bevægelse har nøjagtig samme periode, hvilket er præcis grunden til,
at Kepler-banen er en \emph{lukket} ellipse i stedet for en åben
rosette-bane (jf.\ Bertrands sætning, som siger, at kun
$V\propto1/r$ og $V\propto r^2$ giver lukkede baner for alle
energier). I torus-billedet fra figur~\ref{fig:torus} er
Kepler-banens torus derfor altid \emph{resonant} — den er et
grænsetilfælde af degeneration, ikke blot et specialtilfælde af et
kommensurabelt frekvensforhold.

\begin{bemaerkning}
Sæt Bohr-Sommerfeld-kvantiseringen $J_r=n_rh$, $J_\phi=n_\phi h$
($n_r=0,1,2,\dots$, $n_\phi=1,2,3,\dots$) ind i \eqref{eq:kepler-E-af-J}:
\begin{equation}
  E_n = -\frac{2\pi^2mk^2}{n^2h^2}, \qquad n\equiv n_r+n_\phi = 1,2,3,\dots
\end{equation}
Dette er, med $k=e^2/(4\pi\varepsilon_0)$, netop Bohrs formel for
hydrogenatomets energiniveauer (kompendiets ældre kvanteteori-kapitel,
kapitel~\ref{kap:adiabatisk}, viser metoden generelt; her ser vi dens
mest berømte anvendelse). At $E_n$ kun afhænger af \emph{summen}
$n=n_r+n_\phi$ — og altså er den samme for mange forskellige
kombinationer $(n_r,n_\phi)$ — er den klassiske forklaring på
hydrogenatomets berømte ``tilfældige'' (dynamiske) degeneration:
tilstande med samme $n$, men forskellig fordeling mellem radial og
azimutal bevægelse, har samme energi. Dette er intet mindre end
\emph{fordi} Kepler-banerne er lukkede, altså fordi torusen er
resonant for alle energier, ikke kun langs en isoleret
delmængde af dem. Den fulde kvantemekaniske forklaring på denne
degeneration (via en skjult $SO(4)$-symmetri for Coulomb-potentialet)
gives i [[kvantemekanik-kompendiet]].
\end{bemaerkning}

\begin{bemaerkning}
Dette eksempel illustrerer også, hvorfor degeneration er en
\emph{ikke-generisk} (skrøbelig) egenskab: den opstår her, fordi
$V\propto1/r$ er et af de meget få potentialer med ekstra, skjult
symmetri. Fjernes denne særlige struktur — fx ved at tilføje en lille
forstyrrelse til potentialet, $V(r)=-k/r+\varepsilon f(r)$ — ophæves
degenerationen typisk, $\nu_r\neq\nu_\phi$ for de fleste energier, og
banen bliver en åben rosette i stedet for en lukket ellipse. Hvad der
sker med de \emph{resonante} tori under en sådan forstyrrelse (som
her ligger tæt op ad hele torus-familien, ikke kun en tynd
delmængde) er netop det spørgsmål, KAM-sætningen og det planlagte
kaos-kompendium besvarer.
\end{bemaerkning}
\end{eksempel}

\section{Opsummering}

\begin{itemize}
  \item For periodiske systemer (libration eller rotation) i én
        frihedsgrad defineres virkningsvariablen $J=\oint p\,dq$
        — arealet (op til $2\pi$), som faserumsbanen indeslutter.
  \item $J=J(E)$; den konjugerede vinkelvariabel $w=\partial W/\partial J$
        vokser lineært med tiden, $w(t)=\nu t+w_0$, med
        $\nu=dE/dJ$.
  \item $\nu$ er præcis den fysiske svingningsfrekvens,
        $\nu=1/\tau$ — beregnet uden at skulle finde $q(t)$ eksplicit.
  \item For den harmoniske oscillator: $E(J)=\nu J$ med
        $\nu=\omega/(2\pi)$ — en form, der peger direkte mod Plancks
        kvantiseringsbetingelse.
  \item For separable systemer i flere frihedsgrader defineres en
        virkningsvariabel $J_i$ og frekvens $\nu_i=\partial E/\partial J_i$
        for hver koordinat, som løber uafhængigt og lineært; geometrisk
        ligger bevægelsen på en $n$-dimensional torus i faserummet.
  \item Er frekvenserne $\nu_i$ kommensurable, lukker banen sig selv
        (\emph{resonans}); er de inkommensurable, fylder den torusen
        tæt op (\emph{kvasiperiodisk}). Resonante tori er dem, der
        typisk opløses under en svag forstyrrelse — KAM-sætningens
        emne.
  \item Kepler-problemet ($V=-k/r$) er fuldt degenereret: $E$ afhænger
        kun af summen $J_r+J_\phi$, hvilket forklarer både banens
        lukkethed og hydrogenatomets ``tilfældige'' kvantedegeneration.
\end{itemize}

\begin{opgave}
En partikel med masse $m$ og energi $E$ bevæger sig frit (``rotation'')
på en cirkel med radius $R$ ($0\leq\phi<2\pi$, periodisk). Hamilton-
funktionen er $H=p_\phi^2/(2mR^2)$. Beregn $J=\oint p_\phi\,d\phi$
(integreret over $\phi\in[0,2\pi)$), find $E(J)$, og bekræft, at
$\nu=dE/dJ$ stemmer med den kendte omløbsfrekvens $v/(2\pi R)$, hvor
$v$ er partiklens (konstante) banehastighed.
\end{opgave}

\begin{opgave}
Vis, ved at differentiere \eqref{eq:def-w} med hensyn til tiden og
bruge kædereglen, at $\dot w = (\partial^2W/\partial q\partial J)\dot q$,
og brug dette sammen med \eqref{eq:hamilton-vinkel} til at genudlede
\eqref{eq:delta-w-1} på en alternativ måde: ved at integrere $\dot w$
over én fuld periode $\tau$ i $q$, og sammenholde med
$\Delta w = \nu\tau$.
\end{opgave}

\begin{opgave}
Udled \eqref{eq:kepler-Jr} fra integralet \eqref{eq:kepler-Jr-integral}.
Vink: substituér $u=1/r$ (så $dr=-du/u^2$), hvorved integranden
bliver kvadratroden af et andengradspolynomium i $u$; kvadratkomplettér
dette polynomium, brug standardresultatet
$\oint_{u_1}^{u_2}\sqrt{(u_2-u)(u-u_1)}\,du = \tfrac{\pi}{2}(u_2-u_1)$
(for \'et fuldt omløb, ud og tilbage), hvor $u_1,u_2$ er
polynomiets rødder — udtryk til sidst $u_1,u_2$ ved $m,E,\ell,k$
og reducér til \eqref{eq:kepler-Jr}. Alternativt: bekræft resultatet
numerisk for et selvvalgt talsæt $m,k,E,\ell$.
\end{opgave}
